\section{Preliminary bored-pile design}
\subsection{Geometry and design approach}
The preferred concept uses nominal \SI{600}{\milli\metre} diameter rotary bored piles with toe
level at approximately \SI{17.5}{\metre} below existing grade. Piles extend through the weak
near-surface soils and develop resistance primarily in the two sand units. The preliminary
compression resistance is evaluated by the static sum
\begin{equation}
Q_u=Q_s+Q_b=\sum_i \pi D L_i f_{s,i}
 + \frac{\pi D^2}{4}q_b,
\label{eq:pile}
\end{equation}
where $f_{s,i}$ is the unit shaft resistance and $q_b$ is the unit base resistance. Values are
deliberately capped to reflect bored-pile construction effects \cite{tomlinson2008}.

\subsection{Axial compression resistance}
\begin{table}[H]
\centering
\caption{Static axial resistance calculation for one \SI{600}{\milli\metre} pile.}
\label{tab:pilecalc}
\small
\begin{tabularx}{\textwidth}{@{}L{36mm}C{23mm}C{22mm}C{25mm}Y@{}}
\toprule
\rowcolor{navy!8}
\textbf{Component} & \textbf{Length / area} & \textbf{Unit resistance} & \textbf{Resistance} & \textbf{Basis} \\
\midrule
Clay shaft & \SI{2.5}{m} & \SI{22}{kPa} & \SI{104}{kN} & Adhesion capped \\
Medium-dense sand shaft & \SI{6.0}{m} & \SI{45}{kPa} & \SI{509}{kN} & Effective-stress screen \\
Dense-sand shaft & \SI{7.0}{m} & \SI{80}{kPa} & \SI{1056}{kN} & Bored-pile cap \\
Dense-sand base & \SI{0.283}{m^2} & \SI{3500}{kPa} & \SI{990}{kN} & Clean, sound base \\
\midrule
\multicolumn{3}{r}{\textbf{Calculated ultimate resistance, $Q_u$}} & \textbf{\SI{2659}{kN}} & \\
\multicolumn{3}{r}{Global factor of safety} & 2.5 & \\
\multicolumn{3}{r}{Calculated allowable resistance} & \SI{1064}{kN} & \\
\multicolumn{3}{r}{\textbf{Adopted preliminary allowable load}} & \textbf{\SI{900}{kN}} & test required \\
\bottomrule
\end{tabularx}
\end{table}

The adopted load is less than the static calculation to provide an initial allowance for layer
variability, group interaction, construction tolerance and settlement. It is not a substitute for
code-compliant resistance factors or load testing. Structural axial capacity, reinforcement,
durability, handling of accidental eccentricity and pile-cap connection shall be checked
separately.

\subsection{Pile groups and settlement}
Pile groups should use a preliminary center-to-center spacing of at least $3D$ (\SI{1.8}{\metre})
and preferably $3.5D$ where pile-cap geometry allows. Each group shall be checked for:
\begin{itemize}
  \item sum of individual-pile resistances and block failure;
  \item group settlement under sustained service load, including consolidation of clay below
        and around the equivalent raft;
  \item load redistribution from pile-position tolerance and cap stiffness;
  \item downdrag from new fill, groundwater lowering or settling upper clay; and
  \item minimum edge distance, cap punching shear and reinforcement anchorage.
\end{itemize}

For concept design, expect single-pile service settlement of approximately
\SIrange{10}{18}{\milli\metre} and pile-group settlement of
\SIrange{20}{35}{\milli\metre}. These ranges shall be replaced by a calibrated numerical or
equivalent-raft analysis after the preliminary load test.

\subsection{Uplift and lateral loading}
Do not apply the compression value to uplift. For preliminary proportioning only, limit allowable
tension to \SI{400}{\kilo\newton} per pile, ignoring base resistance and checking the weight of
the pile-cap/structure. Lateral resistance depends strongly on pile-head fixity, group spacing,
cyclic degradation and the weak upper clay. A site-specific $p$--$y$ or continuum analysis is
required once shear, moment and head-displacement criteria are issued by the structural engineer.

\begin{keybox}[Preliminary pile schedule]
\begin{tabularx}{\linewidth}{@{}L{42mm}Y@{}}
Nominal diameter & \SI{600}{\milli\metre} \\
Toe level & \SI{17.5}{\metre} below existing grade, subject to bore log \\
Minimum dense-sand socket & \SI{5.0}{\metre} \\
Allowable compression & \SI{900}{\kilo\newton} per pile \\
Allowable tension & \SI{400}{\kilo\newton} per pile, preliminary only \\
Minimum spacing & $3D$; prefer $3.5D$ \\
Working-load test & \SI{1.5}{\mega\newton} proof load unless governing code requires more \\
Preliminary test pile & Load to geotechnical failure or at least $2.5$ times working load \\
\end{tabularx}
\end{keybox}

% =============================================================================
